\section{Interpreting reals in topological toposes}

We begin interpreting reals with the largest field first.
Fix a space \(X\) to work over.

\subsection{Interpreting MacNeille reals}

\begin{theorem}
  The object of MacNeille reals over \(X\) is generated by tight pairs of semi-continuous real-valued functions,~i.~e.
  \[
    \carr \Rm ≔ \set{\p{\uline f, \oline f} : \uline 𝒞(U, ℝ) × \oline 𝒞(U, ℝ)}{U ∈ 𝒪X, \mathrm{tight}}
  \]
  where \(\mathrm{tight}\) is
  \[
    \uline f(t) = \liminf_{y → t}\oline f(y)\qquad\text{ and}\qquad
    \oline f(t) = \limsup_{y → t}\uline f(y)\text,
  \]
  with equality
  \[
    \eq f g ≔ \int\set{t ∈ U_f ∩ U_g}{\uline f(t) = \uline g(t) ∧ \oline f(t) = \oline g(t)}
  \]
\end{theorem}
\begin{proof}
  First we show that lower (resp. upper) cuts correspond to lower (resp. upper) semi-continuous real-valued functions.

  Let \(L ∈ 𝒫ℚ\) be a lower cut.
  Let \(U\) be the extent of \(L\).
  Then construct the map (in the meta-theory)
  \[ \mordef{\hat L}{U}{𝒫ℚ}{t}{\set{q ∈ ℚ}{t ∈ \i{\uline q ∈ L}}\text.} \]
  The images \(f(t)\) are lower cuts.
  \begin{itemize}
  \item Since \(L\) is inhabited (by \(q : ℚ\)) we have \(q(t) ∈ \hat L(t)\).
  \item Let \(q < p\) and \(q ∈ \hat L(t)\).
    Then internally \(\uline p ∈ L\), so \(p ∈ \hat L\).
  \item Let \(p ∈ \hat L(t)\).
    Then internally we have some \(q > \uline p\) such that \(q ∈ L\).
    It follows that \(q(t) ∈ \hat L(t)\).
  \end{itemize}

  Thus we can view the map \(\hat L\) to be of type \(U → ℝ\).
  Now let \(\p{p, ∞}\) be a basic open.
  It's preimage under \(\hat L\) is \(\set{t ∈ U}{t ∈ \i{\uline p ∈ L}}\).
  As \(\i{p ∈ L}\) is open so is \(\hat L⁻¹\p{p, ∞}\) and \(\hat L\) is lower semi-continuous.

  Conversely, if \(f : U → ℝ\) is lower semi-continuous we construct
  \[ L(t) ≔ \set{q : ℚ}{q(t) < f(t)}\text. \]
  This is a lower cut internally at each \(t ∈ U\).
  \begin{itemize}
  \item Since \(f(t)\) is inhabited (by \(q ∈ ℚ\)) we have \(\uline q ∈ L(t)\).
  \item Let \(q < p\) and \(q ∈ L(t)\). Then \(q(t) < p(t) ∈ f(t)\) so \(q(t) ∈ f(t)\).
  \item Let \(p ∈ L(t)\).
    Then for some \(q > p(t)\) we have \(q ∈ f(t)\).
    Thus \(\uline q ∈ L(t)\).
  \end{itemize}
  Because the property of being a lower cut is geometric we have that \(L\) is a lower cut.

  We have shown that MacNeille reals correspond to \emph{some} pairs of semi-continuous functions.
  The properties~\ref{real-m:L} and~\ref{real-m:U} correspond to tightness.
  TODO: analysis?
\end{proof}


\subsection{Interpreting Dedekind reals}

\begin{proposition}
  For \(f ∈ \carr \Rm\) we have \(\uline f = \oline f\) iff internally the MacNeille real is located.
\end{proposition}
\begin{proof}
  Suppose first the real \(x : \Rm\) is located.
  Write \(L\) for the map into lower cuts corresponding to \(\uline f\) and \(U\) for the map into upper cuts corresponding to \(\oline f\).
  Then we have externally for all \(t ∈ \e x\) and \(p < q\) that \(p < \uline f(t)\) or \(q > \oline f(t)\).
  In the limit towards \(\uline f(t)\) this gives that \(\uline f(t) = \oline f(t)\).

  Conversely, let \(f ∈ \carr \Rm\) be such that \(\uline f = \oline f\) and take \(p < q\) internally.
  Then at \(t ∈ \e f\) we have \(p(t) < f(t)\) or \(f(t) < q(t)\).
  As both of these propositions are open, we have \(p < f ∨ f < q\) internally.
\end{proof}

\begin{corollary}
  The object of Dedekind reals over \(X\) is the sub-Heyting valued set of \(\Rm\) generated by
  continuous real-valued functions,~i.~e.
  \[ \carr \Rd ≔ ⋃_{U ∈ 𝒪X}𝒞(U, ℝ)\text. \]
\end{corollary}


\subsection{Interpreting Cauchy reals}

\begin{example}
  Take the generic real over the countable Fort space, i.~e.~the subset \(\{0\} ∪ \set{2⁻ⁿ}{n ∈ ℕ}\) of the Euclidean plane.

  This is a Cauchy real with Cauchy sequence \(\p{0, n} ↦ 0\) and
  \[ \p{2⁻ᵐ, n} ↦ \begin{cases} 2⁻ᵐ; \quad m < n\\0; \quad m ≥ n\text. \end{cases} \]
\end{example}

\begin{proposition}
  The object of Cauchy reals over \(X\) is the sub-Heyting valued set of \(\Rd\) generated by continuous real-valued functions that are constant on connected subsets.
\end{proposition}
\begin{proof}
  Let \(x : \Rc\) be a Cauchy real.
  Taking any Cauchy sequence for \(x\) (call it \(\p{qₙ}\)) we can analyse it to a sequence of locally constant rational functions with domain \(\e x\).
  Then at any connected subset of \(\e x\) all \(qₙ\) are constant, so they converge to the same real externally.
  Thus the function representing \(x\) as a Dedekind real is constant on connected subsets.

  Conversely, taking such a map \(f : U → ℝ\), we can form the internal sequence as follows.
  Take a connected component \(C\) of \(U\).
  Since \(f\) is constant on \(C\) it has a unique value \(x\), which has a corresponding Cauchy sequence \(\p{qₙ}\).
  This defines a
  TODO: This proof doesn't work

  TODO: this might not be true... I need to think about more innovative sequences of rationals.

\end{proof}





%%% Local Variables:
%%% mode: latex
%%% TeX-master: "../main"
%%% End:
